\documentclass[11pt]{article}
\usepackage{amsmath,amssymb,amsthm}
\usepackage[margin=1in]{geometry}
\newtheorem{theorem}{Theorem}
\newtheorem{lemma}{Lemma}
\begin{document}
\title{V110: Phase-Compatible Shared-Gate MUX Cycles}
\author{NC0\_k-Avoid Laboratory}
\date{}
\maketitle

A signed essential ternary MUX output has the form
\[
 g=o\oplus \operatorname{MUX}(s\oplus p_s,a\oplus p_a,b\oplus p_b).
\]
Fixing its output gives two binary clauses. Selecting branch $r$ yields an
implication
\[
 x_s=\alpha_r\Longrightarrow x_d=\beta_r,
\]
where the two source phases are opposite and the output target selects the
arrival phase.

\begin{lemma}[Single-shared-output composition]
Suppose two directed branch cycles start at the same selector $v$ with opposite
source phases. Their output-gate sets intersect in exactly one gate $h$. If both
cycles require the same target bit on $h$, then a missing output is constructible.
\end{lemma}

\begin{proof}
On each cycle choose targets so every selected implication arrives at the source
phase of the next selected branch, except that the last implication arrives at
the complement of the first source phase. A cycle beginning at phase $\alpha$
therefore contains
\[
 x_v=\alpha\Longrightarrow x_v=1\oplus\alpha,
\]
so every satisfying assignment must set $x_v=1\oplus\alpha$. The two initial
phases are opposite, hence the two cycles force opposite Boolean values on $v$.
All selected outputs are disjoint except $h$, and by hypothesis the two target
assignments agree on $h$. Thus they combine into one fixed-output contradiction.
Unused MUX clauses only strengthen the formula.
\end{proof}

Start with a V109 one-gate bottleneck for selector $v$, first branch
destinations $d_0,d_1$, and bottleneck output $h$. In the V109 return network,
give every output-gate edge capacity one except $h$, which receives capacity two.
All non-gate internal arcs have capacity two and a super-source sends one unit to
each destination.

\begin{theorem}[V110 upgraded-flow avoider]
If the upgraded network has flow value two, the two decomposed return routes
share exactly $h$, and both lifted cycles require the same target bit on $h$,
then a word outside the circuit range is constructible in deterministic
polynomial time.
\end{theorem}

\begin{proof}
Only $h$ has output-gate capacity two. Therefore an integral two-flow yields two
return paths that are gate-disjoint outside $h$. Prepending the prescribed V109
first branches creates two opposite-phase cycles. The preceding lemma applies.
Max-flow, path decomposition, target comparison, and target lifting are all
polynomial-time operations.
\end{proof}

If the upgraded flow remains one, V110 records a nested bottleneck. If flow two
exists but the two traversals demand different target bits on $h$, V110 records
a phase conflict. Neither case is solved here.

\begin{theorem}[Strict exact-stretch family]
For every $k\ge2$ there is an essential MUX circuit with
\[
 n=4k+2,\qquad m=n+1,
\]
that is Hall-minimal, has two distinct-support central MUX outputs, has no V108
SCC-separated certificate under any ignored-output set, forces V109 into a
single shared bottleneck, and is solved by the V110 shared-gate certificate.
Its V102 strong-affine backdoor is exactly
\[
 \beta=2+4\lceil k/2\rceil=\Theta(n).
\]
\end{theorem}

\begin{proof}[Construction and proof sketch]
Use variables $v$, pre-lobes $A_1,\ldots,A_k$ and $B_1,\ldots,B_k$, hub $w$,
and post-lobes $C_1,\ldots,C_k$ and $D_1,\ldots,D_k$. Add central gates
$\operatorname{MUX}(v,A_1,B_1)$ and $\operatorname{MUX}(v,A_2,B_2)$. Each
pre-lobe forms a directed local cycle and also exposes $w$ as its only exit.
Add the shared gate $h=\operatorname{MUX}(w,C_1,D_1)$. Each post-lobe forms the
analogous cycle with exit to $v$.

There are $2+2k+1+2k=4k+3=n+1$ outputs. Since $v$ is the only repeated selector
and every return from either central destination to $v$ must pass through $h$,
V109 returns the bottleneck $h$. Raising only $h$ to capacity two allows one
route through the $C$ post-lobe and one through the $D$ post-lobe; the two
routes share only $h$ and require the same target there.

For the backdoor upper bound, select both hubs $v,w$ and a minimum vertex cover
of each of the four $k$-cycles, giving $2+4\lceil k/2\rceil$. If $v$ is omitted,
all post-lobe selectors must be selected and the two distinct central supports
force additional pre variables. If $w$ is omitted, all pre-lobe selectors must
be selected and the shared output plus the post-lobe constraints impose no
smaller solution. Hence the displayed value is exact.

Hall minimality follows from explicit alternating-path repairs after deleting
any one output. For V108 separation, deleting a unique-selector lobe output makes
its selector acyclic, deleting $h$ makes $w$ acyclic, while deleting a central
bridge either leaves the terminal in the same cyclic component or requires
breaking a unique lobe output and thereby destroys that lobe cycle.
\end{proof}

\paragraph{Boundary.}
V110 does not solve phase conflicts or nested bottlenecks. It therefore does not
prove all essential MUX/bijunctive circuits are in P, unrestricted
$NC^0_3$-Avoid, a new general circuit lower bound, or P versus NP. Novelty and
priority are not established.

\end{document}
